% %%%

%%%   Самарский государственный технический университет

%%%   Кафедра прикладной математики и информатики

%%%

%%%   Преамбула для статьи в журнал "Вестник Самарского государственного

%%%   технического университета. Серия: «Физико-математические науки»"

%%%   Ориентирована под MikTEx 2.4 for Win32

%%%

%%%   Хотя сам журнал собирается под GNU/Linux на дистрибутиве TeXLive



\documentclass[11pt,a4paper,twoside]{article}



\usepackage[cp1251]{inputenc}

\usepackage[T2A]{fontenc}

\usepackage[english,russian]{babel}

\usepackage{samgtubib}

\usepackage[dvips]{graphicx}

\usepackage[multiple]{footmisc}



\usepackage{samgtu}

\renewcommand{\VestnikYear}{2009} % Год издания Вестника

\renewcommand{\VestnikNumber}{2\,(19)} % Номер Вестника

\renewcommand{\VestnikMonth}{Ноябрь} % Месяц выхода Вестника

\renewcommand{\VestnikData}{01.11.09} % Дата подписания Вестника в печать

\renewcommand{\VestnikUPL}{26,64} % Усл. печ. л.

\renewcommand{\VestnikUIL}{25,73} % Уч.-изд. л.

\renewcommand{\VestnikRegNumber}{479/09} % Регистрационный номер

\renewcommand{\VestnikOrder}{994} % Номер заказа








%


% \begin{document}








\renewcommand{\firstpage}{18}


\renewcommand{\lastpage}{25}


%\Partition{Дифференциальные уравнения и краевые задачи}{Writer's Guidelines} % Раздел Вестника, в который подаётся статья





\msc{76L05, 35Q35}


% Номер УДК


\udk{533.601.1}%


\title{О движении газа за фронтом сильной ударной волны, форма которой близка к некоторой кривой}% Полный заголовок


{О движении газа за фронтом сильной ударной волны \dots}% Заголовок, который идёт в верхний колонтитул





\author{В.\,И.~Богатко\inst{1}, Г.\,А.~Колтон\inst{2}, Е.\,А.~Потехина\inst{1}}% Автор(ы) для заголовка, если авторов несколько указать \inst


{Б\,о\,г\,а\,т\,к\,о~В.\,И., К\,о\,л\,т\,о\,н~Г.\,А.,\ П\,о\,т\,е\,х\,и\,н\,а~Е.\,А.}% Автор(ы) для оглавления и колонтитула через





\institute{Санкт-Петербургский государственный университет,\\


199034, г. Санкт-Петербург, Университетская наб., 7-9. \and


Санкт-Петербургский государственный горный институт,\\


199106, г. Санкт-Петербург, Васильевский остров, 21 линия, 2.}





\email{E-mail}{eap225@gmail.com} %E-mail's авторов





\otherinf{ Богатко Всеволод Иванович "--- старший научный сотрудник; к.ф.-м.н., с.н.с. \\


Колтон Гарри Абрамович "--- доцент;  к.ф.-м.н., доцент. \\


 Потехина Елена Александровна "--- старший научный сотрудник; к.ф.-м.н.}





\keywords{газовая динамика, ударная волна, нелинейные уравнения,


уравнения в частных производных}





\workabstract{В работе рассматривается плоская автомодельная


задача о движении невязкого газа за фронтом интенсивной ударной


волны. Предполагается, что форма ударной волны близка к некоторой


кривой, форма которой известна. Решение строится в виде рядов по


степеням малого параметра, характеризующего отношение плотностей


газа на фронте ударной волны. Рассматриваются случаи, когда форма


сильной ударной волны мало отличается от прямой или от окружности.


Решение задачи сводится к интегрированию уравнения


Эйлера"--~Дарбу.}





\engtitle{On Gas Flow Beyond Strong Shock Wave Front, Form of Which Approaches a Certain Curve}





\engauthor{V.\,I.~Bogatko\inst{1}, G.\,A.~Kolton\inst{2},


E.\,A.~Potekhina\inst{1}}{B\,o\,g\,a\,t\,k\,o~V.\,I., K\,o\,l\,t\,o\,n~G.\,A., P\,o\,t\,e\,k\,h\,i\,n\,a~E.\,A.}





\enginstitute{Saint-Petersburg State University,\\


7-9, Universitetskaya nab., St. Petersburg, 199034.


\and


Saint-Petersburg State Mining Institute, \\


2, 21st line, Vasilivcky ostrov, St. Petersburg, 199106.


}% Место работы каждого из авторов на англ. языке








\engabstract{The plane auto model problem of the in viscid gas motion beyond intensive shock wave is studied. It is supposed, that shock wave front approaches some curve, the form of which is known. Solution is constructed in the form of series on the small parameter degrees. This parameter characterizes the relation of gas densities at shock wave front. Certain cases are studied as examples: when intensive shock wave front form is closely approximated to the straight line or to the circle. Solution of the problem is reduced to the Euler--Darboux equation integration.}





\engkeywords{gas dynamics, shock wave, nonlinear equations, partial differential equations}





\date{30/VII/2008} % Дата поступления статьи в редакцию.


\revisiondate{20/II/2009} % В окончательном виде





\otherenginfm{Bogatko Vsevolod Ivanovich, Ph. D. $($Phys. \& Math.$)$.\\


Kolton Garry Abramovich, Ph. D. $($Phys. \& Math.$)$, Ass. Prof.\\


Potekhina Elena Alexandrovna, Ph. D. $($Phys. \& Math.$)$.


}








%%%%%%%%%%%%%%%%%%%%%%%%%%%%%%%%%%%%%%%%%%%%%%%%%%%%%%%%%%%%%%%%%%%%%%%%%%%%%%%%%%%%








\maketitle





Ударные волны представляют собой одно из наиболее важных явлений,


встречающихся в газовой динамике и аэродинамике больших скоростей,


так как они присутствуют во многих течениях, имеющих практическое


значение.











Рассмотрим плоскую задачу о движении невязкого газа за фронтом


ударной волны. Система уравнений газовой динамики, описывающих


течение газа за фронтом волны, может быть записана в виде


\cite{Lunev}:


\begin{equation}


\frac{d\vec v}{dt}=- \frac{1}{\varrho} {\rm grad \,}  p, \quad


\frac{d\varrho}{dt} + \varrho {\rm div \,} \vec v = 0,\quad


\frac{di}{dt} = \frac{1}{\varrho} \frac{dp}{dt}.


\end{equation}


Здесь $p$ "--- давление, $\varrho$ "--- плотность, $\vec v$ --


вектор скорости частиц газа за фронтом ударной волны, $i$ --


энтальпия, $t$ "--- время.





Для того  чтобы система уравнений была замкнутой к ней необходимо


добавить уравнение состояния газа $f(p, \varrho, i)=0$.





Граничными условиями для системы уравнений (1) являются условия


динамической совместности на фронте ударной волны:


\begin{eqnarray}


\varrho \left(\vec v -\vec N_F\right) \vec


n_F&=&\varrho_\infty\left(\vec


v_\infty-\vec N_F\right) \vec n_F, \nonumber\\


(p-p_\infty) \vec n_F&=&\varrho_\infty\left[\left(\vec


v_\infty-\vec N_F\right) \vec n_F\right] (\vec


v_\infty-\vec v), \nonumber\\


p (\vec v  \vec n_F)-p_\infty(\vec


v_\infty\vec n_F)&=&\varrho_\infty\left(\vec


v_\infty-\vec N_F\right)\vec


n_F \left(e_\infty-e+\frac{v_\infty^2}{2}-\frac{v^2}{2}\right), \nonumber


\end{eqnarray}


где индексом $\infty$ обозначены параметры газа перед фронтом


ударной волны,  $\vec n_F$,  $\vec N_F$


"--- орт нормали и скорость перемещения поверхности ударной волны


соответственно,  $e$ "--- внутренняя энергия газа.





Будем рассматривать достаточно сильные ударные волны,  скорость


перемещения которых превышает 3000 м/сек. Как показывают расчеты,


влияние реальных свойств газа на газодинамические параметры потока


за фронтом ударной волны достаточно хорошо можно учесть с помощью


эффективного показателя адиабаты $\gamma$\cite{Lunev}. Тогда


уравнение состояния можно взять в квазисовершенном виде


$$ \frac{1}{\varrho} \equiv \omega(i,  p) = {\gamma - 1\over\gamma} {i\over p},  $$


и для внутренней энергии  будем иметь


$$ e=\frac{1}{\gamma - 1} \frac{p}{\varrho}.$$


Здесь $\gamma$ "--- эффективный показатель адиабаты.





Теперь выписанные условия динамической совместности могут быть


преобразованы к виду:


\begin{eqnarray}


\vec v =\vec v_\infty


-(1-\frac{\varrho_\infty}{\varrho})\bigl[(\vec


v_\infty - \vec N_F) \vec n_F\bigr] \vec n_F,  \nonumber\\


 p=p_\infty+(1-\frac{\varrho_\infty}{\varrho}) \varrho_\infty \bigl[(\vec v_\infty -\vec N_F)


\vec n_F\bigr]^2,  \\


{\frac{p}{p_\infty} = \frac{(\gamma+1) \varrho-(\gamma-1) \varrho_\infty}{


(\gamma+1) \varrho_\infty-(\gamma-1) \varrho}}.  \quad \nonumber


\end{eqnarray}





Из аналитических методов,  применяемых в газовой динамике для


решения задач с сильными ударными волнами,  наиболее широко


используется метод тонкого ударного слоя (метод <<пограничного


слоя>> Г.\,Г.~Черного)~\cite{Chernii}.





Этот метод является одним из вариантов метода малого параметра.


Не\-смотря на то,  что сам малый параметр отсутствует в исходной


системе уравнений газовой динамики в явном виде,  он легко


усматривается в граничных условиях задачи. Метод Г.\,Г.~Черного


основан на естественном предположении о малости (по сравнению с


единицей) отношения плотностей газа перед фронтом сильной ударной


волны и непосредственно за ней. При этом газодинамические


параметры течения представляются в виде рядов специального вида по


степеням параметра ,  характеризующего отношение плотностей газа


перед волной к плотности газа непосредственно за ней.





В настоящее время нет доказательства сходимости,  а значит и


строгого математического обоснования,  метода тонкого ударного


слоя. Поэтому для оценки практической сходимости этого метода было


построено третье приближение для определения параметров газа в


ударном слое клина и конуса при движении с большой переменной


скоростью. Сравнение второго и третьего приближений показало,  что


с достаточной степенью точности в расчетах можно ограничиться


двумя приближениями \cite{Potekhina}.





Расчеты показывают,  что хорошее совпадение приближенных решений с


точными решениями и численными исследованиями может быть получено,


если в разложении искомых функций в ряд удерживать не менее двух


членов (cм.,  например,  [\citen{Chernii1,  Zapryanov}]).





В настоящей работе будем рассматривать автомодельные течения газа


за фронтом интенсивной ударной волны.





Пусть сильная ударная волна движется по покоящемуся газу


($v_\infty = 0$). При построении аналитических решений задач с


сильными ударными волнами обычно в условиях на фронте ударной


волны пренебрегают противодавлением $p_\infty$ и квадратом


отношения скорости звука перед фронтом ударной волны к скорости ее


перемещения \cite{Sedov}.





Будем считать,  что форма фронта сильной ударной волны близка к


некоторой кривой $(\ell)$,  уравнение которой в плоскости


безразмерных автомодельных переменных запишем в виде $r=r(s)$,  где


$s$ "--- координата вдоль этой кривой.





Тогда в системе координат,  связанной с данной кривой, уравнения


газодинамики (1) в безразмерных переменных примут вид:


\begin{gather}


\bigl(v - r\cos\alpha\bigr) \frac{\partial u}{\partial


s} + (u - \lambda +  r\sin\alpha) \frac{\partial


u}{\partial \lambda} + k v \bigl(v - r


\cos\alpha\bigr) = - \omega \frac{\partial p}{\partial


\lambda},\\


\bigl(v-r\cos\alpha\bigr)\frac {\partial v}{\partial


s}+(1-k\lambda)(u-\lambda+r\sin\alpha)\frac {\partial v}{\partial


\lambda}-ku\bigl(v-r\cos\alpha\bigr)=-\omega\frac{\partial


p}{\partial s},\\


\bigl(v-r\cos\alpha\bigr)\frac{\partial\omega}{\partial


s}+(1-k\lambda)(u-\lambda+


r\sin\alpha) \frac{\partial\omega}{\partial\lambda} =\omega \Bigl\{


\frac{\partial v}{\partial s} + {\partial


[(1-k\lambda)u]\over\partial\lambda}\Bigr\},


\end{gather}


\begin{multline}


\bigl(v - r\cos\alpha\bigr) \frac{\partial i}{\partial


s} + (1-k\lambda) (u -


\lambda + r\sin\alpha) \frac{\partial i}{\partial \lambda} =\\


=\omega \Bigl\{\bigl(v - r\cos\alpha\bigr) \frac{\partial


p}{\partial s} +


 (1-k\lambda) (u - \lambda + r \sin\alpha)  \frac


{\partial p}{\partial \lambda} \Bigr\},


\end{multline}


где $u$,  $v$ "--- составляющие вектора скорости по нормали и


касательной к кривой $(\ell)$ соответственно,  $\lambda$ ---


координата по нормали,  $\omega$ "--- удельный объем,  $\alpha$ ---


угол между радиусом-вектором кривой и вектором,  касательным к ней,


$k$ "--- кривизна кривой $(\ell)$,  компоненты вектора скорости


отнесены к скорости звука  $a_\infty$ перед фронтом,  давление ---


к ${\varrho_\infty} {a^2_\infty}$,  плотность "--- к


$\varrho_\infty$,  энтальпия "--- к


 $a^2_\infty$,  время "--- к отношению характерного размера к


 скорости звука перед фронтом. Характерный размер выбирается в


 зависимости от специфики конкретной задачи.





Граничными условиями для системы уравнений (3)--(6) являются


условия динамической совместности на фронте ударной волны,  также


записанные в безразмерном виде,  при этом скорость перемещения


фронта ударной волны отнесена к скорости звука перед фронтом.





Будем решать задачу методом тонкого ударного слоя.  Представим


уравнение фронта ударной волны в виде


\[


\vec{r}_{F}(s) = \vec{r}(s) + \varepsilon


r_1(s) \vec{n},


\]


где $\vec{n}$ "--- орт нормали к кривой $(\ell)$,


$\varepsilon$ "--- малый параметр,  характеризующий отношение


плотностей газа на фронте интенсивной ударной волны,  а индекс $F$


указывает на то,  что величина вычисляется на фронте ударной волны.





Из условий динамической совместности получим


\[


N = \vec{r}_{F}\cdot\vec{n}_{F}=


 N_0 +\varepsilon N_1,


\]


\[


v_F = \varepsilon r_1^\prime N_0,  \quad


u_F = - N_0 - \varepsilon \bigl(N_1 + r_1 k - \omega_0


\bigr),


\]


\[


p_F = p_{0F} + \varepsilon p_{1F},  \quad


i_F = i_{0F} + \varepsilon i_{1F},


\]


где  $  N_0 = r \sin\alpha$; $  N_1 = r r_1  cos


\alpha - k r r_1 \sin\alpha - r_1$,    $p_{0F} = N_0^2$,


 $\omega_0 = \omega_F/\varepsilon$.





В ряде задач газовой динамики с сильными ударными волнами форма


фронта ударной волны на сравнительно небольших и подчас наиболее


интересных для приложений участках мало отличается от окружности


или прямой. При этом характерная длина дуги ударной волны имеет


порядок $\sqrt\varepsilon$. К~таким задачам относятся,  например,


задача дифракции сильной ударной волны около угла,  задача


отражения сильной ударной волны от твердой стенки,  задача о поршне


и др.





Так как основная масса газа в таких задачах сосредоточена в узкой


зоне,  примыкающей к фронту ударной волны,  положим


\[


\lambda = \varepsilon \lambda_0,  \quad


s = \sqrt\varepsilon s_0.


\]





Далее имеем


\[


r =


\begin{cases}


\:\:\:1  & \mbox{(для окружности)}, \\


 \frac{1}{\sin\alpha}


&\mbox{(для прямой)}.


\end{cases}


\]


Тогда получим $\quad N_0 = 1$,  $p_{0F} = 1$,


\[


\cos\alpha = \frac{dr}{ds} \;


\begin{cases} \equiv 0


 & \mbox{(для окружности)},\\


=\sqrt\varepsilon s_0 \bigl(1 - \frac{\varepsilon}{


2} s_0^2\bigr)  + \dots & \mbox{(для прямой)},


\end{cases}


\]


\[


r_1^{\prime} \cos\alpha = \frac{1}{\sqrt\varepsilon} \frac{dr_1}


{ds_0} \cos\alpha =


\begin{cases} \:\:\:0 &\mbox{(для окружности)},\\


 s_0 \frac{dr_1}{ds_0} &\mbox{(для прямой)},


\end{cases}


\]


\[


 k=\begin{cases} 1  &\mbox{(для окружности)},\\


 0 &\mbox{(для прямой)}.


 \end{cases}


\]


 Объединяя эти выражения,  их можно


записать так:


\[


r \cos\alpha = \sqrt\varepsilon h_0,


\]


где


\[


h_0 = \begin{cases} 0  &\mbox{(для окружности)},\\


s_0  &\mbox{(для прямой)}.


\end{cases}


\]


Таким образом,  условия на фронте ударной волны окончательно примут следующий


вид:


\[


N_0 = 1,  \quad N_1 =


\begin{cases}


-2r_1 &\mbox{(для окружности)},\\


s_0 \frac{dr_1}{ds_0} - r_1 &\mbox{(для прямой)},


\end{cases}


\]


\[u_F = u_{0F} + \varepsilon u_{1F},   \]


\[


u_{0F}=-1,  \quad   u_{1F} =


\begin{cases}


2r_1 - r_1 + \omega_0 = r_1 +  \omega_0 &\mbox{(для окружности)},\\


- s_0 \frac{dr_1}{ds_0} + \omega_0  &\mbox{(для прямой)},


\end{cases}


\]


\begin{equation}


v_F =  \sqrt\varepsilon v_{1F},  \quad  v_{1F} =


\frac{dr_1}{ds_0},


\end{equation}


\[


p_{0F} = 1,  \quad i_{0F} =  \frac{1}{2}, \quad


\omega_{0F} = 1 \quad \mbox{и  т.\,д.}


\]





 Исходя из анализа граничных условий,  искомые функции будем искать


 в виде рядов следующего вида:


\[


u=u_0+\varepsilon u_1 +\cdots,  \quad v=\sqrt\varepsilon v_1+


\cdots, \quad p=p_0 + \varepsilon p_1 +\cdots,


\]


\begin{equation}


i= i_0 + \varepsilon i_1 + \cdots, \quad


\omega = \varepsilon


\left(\omega_0 + \varepsilon \omega_1 + \cdots  \right).


\end{equation}





Подставляя (8) в систему уравнений (3)--(6),  получим для главных


членов разложения следующую систему уравнений:


\[


\left(u_0 + 1\right) \frac{\partial u_0}{\partial


\lambda_0} = 0,  \quad \left(u_0 + 1\right) \frac{\partial


v_1}{\partial \lambda_0} = 0,


\]


\[


\left(u_0 + 1\right) \frac{\partial \omega_0}{\partial


\lambda_0} = \omega_0  \frac{\partial u_0}{\partial \lambda_0},


\quad \left(u_0 + 1\right) \frac{\partial i_0}{\partial


\lambda_0} = 0.


\]


Из первого уравнения этой системы с учётом граничных условий (7)


получим


\begin{equation}


u_0 = u_{0F}(s_0) = - 1.


\end{equation}





Уравнение (3) с учётом (9) даёт


\[


\frac{\partial p_0}{\partial \lambda_0} = 0,


\]


откуда очевидно,  что с учётом граничных условий (7) будем иметь


\begin{equation}


p_0 = p_0(s_0) \equiv 1.


\end{equation}





Учитывая в системе уравнений (3)--(6) члены более высокого


порядка,  с учётом (9) и (10) получим систему уравнений для


определения следующих членов разложения:


\begin{gather}


(v_1 - h_0)\frac{\partial u_1}{\partial


s_0}+(u_1-\lambda_0)\frac{\partial u_1}{\partial


\lambda_0}+k v_1 (v_1-h_0) =-\omega_0\frac{\partial


p_1}{\partial \lambda_0},\\


(v_1 - h_0) \frac{\partial v_1}{\partial


s_0} + (u_1 - \lambda_0) \frac{\partial v_1}{\partial


\lambda_0}  +k (v_1-h_0) = - \omega_0 \frac{\partial


p_0}{\partial s_0},\\


(v_1 - h_0) \frac{\partial \omega_0}{\partial


s_0} + (u_1 - \lambda_0)  \frac{\partial \omega_0}{\partial


\lambda_0} =  \omega_0 \left\{\frac{\partial v_1}{\partial


s_0} +  \frac{\partial u_1}{\partial \lambda_0} + k\right\},\\


(v_1 - h_0) \frac{\partial i_0}{\partial


s_0} + (u_1 - \lambda_0)  \frac{\partial i_0}{\partial


\lambda_0} = 0,\\


(v_1 - h_0) \frac{\partial i_1}{\partial


s_0} + (u_1 - \lambda_0)  \frac{\partial i_1}{\partial


\lambda_0} = 0.


\end{gather}


Из уравнения (14) следует,  что \ $i_0$ \ постоянно вдоль


траектории.





Тогда с учётом граничных условий (7) получаем


$ i_0 \equiv  \frac{1}{2}$,  а так как


$p_0 \equiv 1$,


то из уравнения состояния имеем


\begin{equation}


\omega_0 \equiv 1.


\end{equation}





Учитывая (16),  из уравнений (12) и (13) получим систему двух


уравнений для определения функций $v_1$ и $u_1$:


\begin{gather}


(v_1 - h_0) \frac{\partial v_1}{\partial


s_0} + (u_1 - \lambda_0)  \frac{\partial v_1}{\partial


\lambda_0} +k (v_1-h_0) = 0, \\


 \frac{\partial


v_1}{\partial s_0} + \frac{\partial u_1}{\partial


\lambda_0} + k = 0.


\end{gather}


Эту систему уравнений удобно переписать в виде


\begin{gather}


(v_1 - h_0) \frac{\partial {(v_1 + k s_0)}}{\partial


s_0} +  \left(u_1-\lambda_0\right) \frac{\partial


{(v_1 + k s_0)}}{\partial \lambda_0} = 0,


\\


\frac{\partial {(v_1 + k s_0)}}{\partial


s_0} + \frac{\partial u_1}{\partial \lambda_0} = 0.


\end{gather}





Из уравнения (20) следует существование функции


$Q(s_0,  \lambda_0)$ такой, что


\[


\frac{\partial Q}{\partial \lambda_0} = v_1 + k s_0,  \quad


\frac{\partial Q}{\partial s_0} = - u_1.


\]


Тогда из уравнения (19) получим нелинейное дифференциальное


уравнение в частных производных второго порядка для функции


$Q(s_0,  \lambda_0)$:


\begin{equation}


\left(\frac{\partial Q}{\partial


\lambda_0} - k s_0-h_0)\right) \frac{{\partial}^2 Q}{\partial


\lambda_0   \partial s_0} -  \left(\frac {\partial Q}{\partial


s_0} + \lambda_0\right) \frac{{\partial}^2 Q}{{\partial


\lambda_0}^2} = 0.


\end{equation}





Для того  чтобы вместо нелинейного уравнения (19) получить


линейное уравнение в частных производных,  положим \cite{Kamke}:


\[


 Q = \Phi(q,  s_0) + \lambda_0 q, \] где


$q = \frac{\partial Q}{\partial \lambda_0}$.


Тогда будем иметь


\begin{gather}


q=\frac{\partial \Phi}{\partial q} \frac{\partial q}{\partial


\lambda_0} + q +\lambda_0 \frac{\partial q}{\partial


\lambda_0},  \quad \lambda_0 = - \frac{\partial \Phi}{\partial


q},  \quad \frac{{\partial}^2 Q}{\partial \lambda_0 \partial


s_0}= \frac{\partial q}{\partial s_0},\notag \\


\frac{\partial Q}{\partial s_0}= \frac{\partial \Phi}{\partial


q} \frac{\partial q}{\partial s_0} + \frac{\partial


\Phi}{\partial s_0} + \lambda_0 \frac{\partial q}{\partial


s_0} = \frac{\partial \Phi}{\partial s_0},\\


\frac{{\partial}^2 Q}{\partial s_0 \partial


\lambda_0} = \frac{{\partial}^2 \Phi}{\partial s_0 \partial


q} \frac{\partial q}{\partial \lambda_0} =  \frac{{\partial}^2


Q}{\partial \lambda_0^2} \frac{{\partial}^2 \Phi}{\partial


s_0 \partial q}.\notag


\end{gather}


Подставляя (22) в (21),  получим


\begin{equation}


\frac{{\partial}^2 Q}{\partial


\lambda_0^2} \left[(q - k s_0 - h_0) \frac{{\partial}^2


\Phi}{\partial s_0 \partial q} -  \left(\frac{\partial


\Phi}{\partial s_0} - \frac{\partial \Phi}{\partial


q}\right)\right] = 0.


\end{equation}


Из (23) очевидно следует,  что могут иметь место два случая.





В первом случае выполняется соотношение


\[


\frac{{\partial}^2 Q}{\partial \lambda_0^2} \equiv 0.


\]


Тогда


$


\frac{\partial v_1}{\partial \lambda_0} = 0$ и


\begin{equation}


v_1 = f_1(s_0), \quad


u_1 = - \left[k + f_1'(s_0)\right] \lambda_0 + f_2(s_0).


\end{equation}


Подставляя (24) в уравнение (17),  получим:





или


\[ f_1(s_0) =


\begin{cases}


0  & \mbox{(для окружности)},\\


s_0  & \mbox{(для прямой), }


\end{cases}


\]





или


\[


f_1(s_0) =


\begin{cases}


- s_0 + c  & \mbox{(для окружности)}, \\


\:\:\:\: c  & \mbox{(для прямой)}.


\end{cases}


 \]





Во втором случае,  если квадратная скобка в (23) равна нулю,  будем


иметь для функции $\Phi(s_0, q)$ линейное дифференциальное


уравнение в частных производных второго порядка:


\[


(q - k s_0 - h_0) \frac{{\partial}^2 \Phi}{\partial s_0


\partial q} - \left(\frac{\partial \Phi}{\partial


s_0} - \frac{\partial \Phi}{\partial q} \right) = 0.


\]


При этом


\[


 h_0 + k s_0 =


\begin{cases}


s_0  \quad & \mbox{(для окружности)},  \\


s_0 \quad & \mbox{(для прямой)}.


\end{cases}


 \]


Отсюда следует,  что для функции $\Phi$ имеем уравнение


Эйлера"--~Дарбу:


\[


\frac{{\partial}^2 \Phi}{\partial s_0 \partial


q} = \frac{1}{q - s_0}\left(\frac{\partial \Phi}{\partial


s_0} - \frac{\partial \Phi}{\partial q}\right).


\]





Как известно общее решение этого уравнения может быть построено с


двумя произвольными функциями~\cite{Trikomi}. Тогда легко может


быть найдено выражение для $u_1$ и $v_1$,  а затем из уравнения


(11) определяется $p_1$. Коэффициент $i_1$ в разложении для


энтальпии найдётся из уравнения (15),  а из уравнения состояния


определяется $\omega_1$.





Таким образом, в настоящей работе предлагается схема получения


приближенного аналитического решения нелинейной задачи определения


параметров газа за фронтом достаточно сильной ударной волны,  форма


которой известна. Эта схема может быть использована для решения


конкретной краевой задачи газовой динамики. Тогда граничное


условие,  определяющее специфику конкретной задачи (условие


обтекания,  условие на поршне и т.п.),  позволит уточнить форму


фронта ударной волны.





\newpage


\begin{thebibliography}{6}


\RBibitem{Lunev}


\by Лунев~В.\,В.


\book  Гиперзвуковая аэродинамика


\publaddr М.


\publ Машиностроение


\yr 1975


\finalinfo 327~c





\RBibitem{Chernii}


\by Черный~Г.\,Г.


\book  Течения газа с большой сверхзвуковой


скоростью


\publaddr М.


\publ Физматгиз


\yr 1959


\finalinfo 220~c





\RBibitem{Potekhina}


\by Потехина~Е.\,А.


\paper Метод тонкого ударного слоя в задаче о


нестационарном обтекании клина и конуса гипрезвуковым потоком


газа


\jour  Вестн. Ленингр. ун-та


\yr 1977


\issue 7


\pages 117--123





\RBibitem{Chernii1}


\by Черный~Г.\,Г.


\paper  Адиабатические движения совершенного газа с


ударными волнами большой интенсивности


\jour Изв. АН СССР,  ОТН


\yr  1957


\issue 3


\pages 66--81





\RBibitem{Zapryanov}


\by Запрянов~З.\,Д.


\paper Численное исследование сверхзвукового


обтекания плоских и осесимметричных затупленных тел при


$M\to\infty$ и $\gamma\to 1$


\jour Изв. АН СССР. Механика жидкости и газа


\yr 1967


\issue 4


\pages 164--167





\RBibitem{Sedov}


\by Седов~Л.\,И.


\book  Методы подобия и размерности в механике


\publaddr М.


\publ Наука


\yr 1972


\pages 440~c





\RBibitem{Kamke}


\by Камке~Э.


\book  Справочник по дифференциальным уравнениям в~частных производных первого порядка


\publaddr М.


\publ Наука


\yr 1966


\finalinfo 260~c





\RBibitem{Trikomi}


\by Трикоми~Ф.


\book Лекции по уравнениям в частных производных


\publaddr М.


\publ Иностр. лит-ра


\yr  1957


\finalinfo 443~c


\end{thebibliography}





%\newpage











\makeengtitlem








  \lfoot[\thepage]{}


  \rfoot[]{\thepage}


%


% \newpage


%


% \tableofengcontents


%


% \end{document}


